\chapter{The report}\label{ch:report}

\section{What a report is}

A report is one specimen's reading and everything needed to check it. Chain v4 writes it as one self-contained HTML page rendered from the run's
bundle. A banner at the top states the chain version, the run, and that it is a research instrument, not a diagnostic test; it also says that each
reading carries its commissioning status, and that a reading outside what an instrument was tested on is shown with a note saying so.

\begin{center}\small
\begin{longtable}{@{}p{0.24\textwidth}p{0.68\textwidth}@{}}
\caption{Sections of the chain v4 report.}\label{tab:tabs}\\
\toprule section & what it holds \\\midrule\endfirsthead
\toprule section & what it holds \\\midrule\endhead
at a glance & one line per instrument: the reading, its state, and its commissioning status or the plain note \\
cell shares & the IAM-Deconvolver's shares against IAM-Atlas v2.1, each with its interval, the commissioned cells marked \\
Met-A & the gauge with Normal marked, $A$ (tared when references were run beside it), its state or the reason it was withheld \\
C-score & $C$, its one-sided band (above $C_{\rm rel}=1.152$ is a clustered departure), and the HEALPix sky plate of the departures \\
IAM-A & when sequencing input is given: the sample-shape check, the held error, the loose share, the held error the line predicts, IAM-A on the gauge with its own ends, and its state \\
the globe & the sample on the atlas's sphere, a reader can turn it: the sample, its healthy reference (the atlas, weighted by the sample's own cell shares), the fixed change (each site's difference in its healthy width and the sample's own spread, read in blocks of 10 neighboring sites, as the C-Score reads them; a block has changed beyond the line pure chance crosses less than once in 20 genomes, so a healthy sample stays nearly dark and a change that runs along the genome lights up in clumps) with a slider that only fades the healthy background, the change at Met-A's 6,000 sites against its purified-array reference, in the same blocks, that reference itself, and each cell the IAM-Deconvolver read with its share; a picture, nothing is scored from it; the same globe is also written on its own full page beside the report \\
technical record & the full stage-by-stage record: intake, calibration, composition, tare, entropy ceiling, refusals, what the build withholds, and the bundle \\\bottomrule
\end{longtable}
\end{center}

The gauge marker is drawn from $A_{\rm rel}$ whenever the tare produced one. Untared isolated neutrophils are drawn against their healthy reference and labeled untared.
Untared whole blood has no gauge position; the number is printed and the reader is told to supply references.

% COMMISSIONING-RETURN fig:p4_reportgauge: figure taken out 2026-10-03 (plotted development values from the removed fitted tare). Put back, redrawn from the commissioned chain; see development/COMMISSIONING_RETURNS.md.


\section{What a report does not say}

The report does not name a condition, and it states its own state words only: \emph{Normal}, \emph{above Normal}, \emph{below Normal}, \emph{untared}. Covariates
supplied with a run are kept in the custody record, the bundle and the ledger, and no stage reads them. Because the bundle block of the
page shows the bundle, a covariate supplied as free text appears there. The v4 page carries a render-time vocabulary guard: the page is not written if its front prose carries a condition name or a word the report may not print.

% COMMISSIONING-RETURN fig:p4_states: figure taken out 2026-10-03 (plotted development values from the removed fitted tare). Put back, redrawn from the commissioned chain; see development/COMMISSIONING_RETURNS.md.


\section{What a reading outside Normal does not mean}\label{sec:p4notmean}

A reading above Normal is a statement about the error of one cell type on its own identity sites, and no more.
\begin{itemize}
\item It does not locate anything: it says that cells of one type with more error than healthy are in the specimen, not where they came
      from or whether they have formed anything.
\item It does not name a cause: a direction on the gauge is a regime, not a mechanism.
\item It is not a grade of a lesion and not a measure of its size. Fidelity of a pattern may track grade in tissue; it is its own axis.
\item It is not a clock: age is carried on the record and is not an operand of any reading.
\item It is not a fuel gauge: the reading is of the pattern, not of the energy spent holding it (Chapter~\ref{ch:landauer}).
\end{itemize}
\derived{} for what the arithmetic contains; the rest is the rule of the report.
The honest reading of a reading above Normal is a prompt for a clinician, not a verdict for a patient: the instrument flags and refers; it
does not diagnose. Chapter~\ref{ch:astrogenetics} sets out in plain language what such a reading does and does not mean, and why a benign
growth is expected to read near 1 (Section~\ref{sec:ag_benign}).

Because the gauge reads identity fidelity and not cell number, two expectations follow. A benign growth whose cells remain well
differentiated should read near 1 on its own cells: those cells have lost growth control, not identity. And a reading does not depend
on a lesion's shape, since it reads DNA rather than an image. \prediction{} Neither is tested. An ordered tissue series (normal, benign, dysplastic, established) read cell by cell is a named test in the plan's precancer track. \openprob

\section{Refusals are printed}

Everything the chain declines to do appears with the reason: a quarantined array (no page; the run stops), a refused platform, a composition not solved, a
neutrophil fraction below the read line, too few identity sites measured, a reading left untared, a methylated-site mean past the entropy ceiling, and an IAM-A declined by the sample-shape check or for too few qualifying molecules. The rule is that a reader should never have to infer what the instrument did not measure.

\section{The report as a laboratory record}

Each run writes, beside the page, the bundle as JSON and one row in an evidence ledger: run number, specimen, platform and array type, intake verdict, floor and
reference versions, neutrophil fraction, $A$, state, $A_{\rm rel}$, tare method and number of references, noise index, detection limit, C-score and, when present,
IAM-A with its pipeline. The intake record in the bundle carries the integrity hash of the IDAT pair. A reading is reproducible if a second laboratory with the same
IDATs, the same references and the same frozen files obtains the same numbers. The bundle carries the chain's commit, and the release check verifies the checksum of every frozen file the chain reads (FROZEN\_INPUTS\_v4.json) on a fresh copy of the repository.

Two readings are comparable only if they were made with the same instrument. The record therefore needs a fingerprint: the chain's
commit and the checksum of every frozen file the run read (the reference, the identity sites, the composition profiles, the noise
sites) and of the tare inputs. Two readings with different fingerprints are compared only after the difference is named, and a reading whose recorded checksums no longer match the repository is not quoted until the difference is resolved. The commit is recorded with every run; writing each frozen file's checksum into every run's own bundle is not yet built. \openprob
